\section*{摘要}
\addcontentsline{toc}{section}{摘要}

IronBuddy 是一个面向居家健身场景的嵌入式智能教练系统。项目关注的核心
问题是：摄像头能够看到动作姿态，却难以判断目标肌肉是否真正参与发力。
在深蹲、弯举这类训练中，两次动作的骨架轨迹可能接近，但目标肌激活、
代偿肌参与和训练刺激并不相同。系统因此把视觉姿态估计与表面肌电
（sEMG）放入同一条反馈链路，尝试同时判断动作形态、目标肌参与和代偿
趋势。

系统由 RK3399ProX 边缘计算节点、摄像头和麦克风等板端外设、ESP32 肌电
传输单元、贴片电极与穿戴固定结构、FSM/GRU 动作评估模块、语音交互、
Web 前端和后台记录接口组成。视觉侧提供关节角度、角速度和动作相位，
肌电侧通过 RMS 与 MVC 归一化描述目标肌和代偿肌的激活水平；两类信号
经过 Mid-Fusion 形成 7D 特征窗口，再由 CompensationGRU 给出标准、
代偿或不标准的动作质量判断。训练反馈不只统计完成次数，也围绕目标肌
激活、代偿情况和疲劳积分组织，更接近结果导向的训练记录方式。

在交互层面，系统通过语音唤醒、指令控制、DeepSeek 对话、训练总结和
本地数据库记录，减少运动过程中对触屏操作的依赖。长期记录、用户偏好
和周期性总结作为工程接口预留，属于后续扩展方向；真实 sEMG 连续输入、
语音现场稳定性与完整手测流程仍需继续复测。

\textbf{关键词}：嵌入式系统、边缘计算、表面肌电、人体姿态估计、多模态
融合、人机交互

\newpage
